\section{Results and Discussion}

TODO: Replace every placeholder below with measured values from the final evaluation. Do not submit the table with placeholders or infer results from training accuracy alone.

\begin{table}[t]
    \caption{Performance on the held-out test set. Replace placeholders with measured values.}
    \label{tab:results}
    \centering
    \begin{tabular}{lcccc}
        \toprule
        Method & Accuracy & Precision & Recall & F1-score \\
        \midrule
        Proposed CNN & -- & -- & -- & -- \\
        Baseline & -- & -- & -- & -- \\
        \bottomrule
    \end{tabular}
\end{table}

Discuss the results in Table~\ref{tab:results} with the dataset split and evaluation unit in mind. TODO: Explain the main findings, compare with the stated baselines, analyze errors or failure cases, and distinguish in-domain performance from cross-dataset generalization. Fig.~\ref{fig:results} is optional.

\begin{figure}[t]
    \centering
    \IfFileExists{figures/results.png}{
        \includegraphics[width=\linewidth]{results.png}
    }{\fbox{\parbox{0.9\linewidth}{\centering Add a results plot as\\\texttt{figures/results.png}}}}
    \caption{Measured evaluation results. Replace this placeholder with a plot generated from the completed experiments.}
    \label{fig:results}
\end{figure}

\subsection{Limitations}
TODO: Discuss dataset coverage, compression and image-quality sensitivity, domain shift, possible demographic effects, and any constraints of the experimental setup that were actually assessed or remain untested.
